\documentclass[a4paper,10pt]{article}

\usepackage[a4paper, top=1.15cm, bottom=1.15cm, left=1.35cm, right=1.35cm]{geometry}
\usepackage{fontspec}
\usepackage{enumitem}
\usepackage{xcolor}
\usepackage[colorlinks=true, urlcolor=navyblue, linkcolor=navyblue]{hyperref}

\definecolor{navyblue}{RGB}{15, 62, 130}

\setmainfont{verdana.ttf}[
    Path = ./fonts/,
    BoldFont = verdanab.ttf,
    ItalicFont = verdanai.ttf,
    BoldItalicFont = verdanaz.ttf
]

\pagestyle{empty}
\setlength{\parindent}{0pt}
\setlength{\parskip}{0pt}

\setlist[itemize]{
    leftmargin=13pt,
    label=\textbullet,
    itemsep=1pt,
    parsep=0pt,
    topsep=1.5pt,
    partopsep=0pt
}

\newcommand{\cvsection}[1]{%
    \vspace{5pt}%
    {\fontsize{10.5pt}{12pt}\selectfont \textbf{#1}}\\[-4pt]%
    \rule{\textwidth}{0.5pt}%
    \vspace{2.5pt}%
}

\begin{document}

\begin{center}
    {\fontsize{15.5pt}{18pt}\selectfont \textbf{Rizky Yanuar Kristianto}} \\[3pt]
    Surabaya, East Java, Indonesia \\
    Phone: \href{tel:+6289676257683}{+6289676257683} \;\textbar\;
    Email: \href{mailto:rizkyyanuarkristianto@gmail.com}{rizkyyanuarkristianto@gmail.com} \\[2pt]
    \href{https://www.linkedin.com/in/rizkyyanuark}{linkedin.com/in/rizkyyanuark} \;\textbar\;
    \href{https://github.com/rizkyyanuark}{github.com/rizkyyanuark}
\end{center}

\vspace{3pt}
Data Science student with hands-on experience in data engineering and applied machine learning. Experienced in building automated ETL pipelines with Apache Airflow and PySpark, working with vector databases for RAG workflows, and deploying containerized applications on AWS and GCP. Focused on building clean data pipelines that support reliable machine learning models.

\cvsection{EDUCATION}

\textbf{STATE UNIVERSITY OF SURABAYA} \hfill Surabaya, East Java \\
\textbf{Bachelor of Data Science} \hfill Aug 2022 -- Jul 2026 \\
\textbf{GPA:} 3.59/4.00 (145 Credits completed) \\
Relevant Coursework: Data Science, Data Structures and Algorithms, Statistics, Data Modeling, Machine Learning, Computer Vision, Generative AI.

\cvsection{WORK EXPERIENCE}

\textbf{State University of Surabaya} \hfill Surabaya, East Java \\
Research Assistant \hfill Jan 2026 -- Present
\begin{itemize}
    \item Co-authored a research paper currently under peer review in a Scopus Q1 journal, focusing on generative model architectures and knowledge graphs.
    \item Ran hyperparameter tuning and ablation experiments to evaluate architectural improvements against baseline models for an IEEE-format manuscript.
    \item Set up Docker environments on AWS GPU instances to run data extraction, ontology mapping, and embedding generation workflows.
    \item Tracked datasets and model checkpoints using Git and DVC to keep experimental results organized and reproducible.
\end{itemize}

\vspace{4pt}
\textbf{Kominfo Jatim} \hfill Surabaya, East Java \\
AI Engineer Intern \hfill Jul 2025 -- Dec 2025
\begin{itemize}
    \item Built a RAG-based question-answering bot for the Majadigi portal using AWS Bedrock and Qdrant, helping citizens search regional regulations easily.
    \item Trained an NLP classification model for hoax news detection, helping the internal team review citizen incident reports more quickly.
    \item Set up GitLab CI/CD pipelines to run automated tests and deploy code to a Linux VPS, and documented REST APIs in Postman for the frontend team.
    \item Assisted in deploying the AI features to production on Majadigi, serving live citizen inquiries.
\end{itemize}

\vspace{4pt}
\textbf{PT PLN Icon Plus} \hfill Surabaya, East Java \\
Data Engineer Intern \hfill Feb 2025 -- Jun 2025
\begin{itemize}
    \item Built automated ETL pipelines using Apache Airflow, PySpark, and Python to process 240,000+ daily ICONNET telecom records, reducing manual reporting time by 50\%.
    \item Created a Slack alert notification linked to Airflow DAGs, enabling the team to detect and resolve pipeline errors quickly.
    \item Analyzed telecom signal quality data across providers and built interactive Looker Studio dashboards for weekly operational monitoring.
    \item Built a database question-answering tool using an LLM and SQL queries, helping 15+ staff look up network records without writing manual SQL.
\end{itemize}

\vspace{4pt}
\textbf{Bangkit Academy led by Google, GoTo, and Traveloka} \hfill Remote \\
Machine Learning Cohort \hfill Aug 2024 -- Dec 2024
\begin{itemize}
    \item Led a 7-member capstone team (ML, Cloud, Mobile) to develop NV-Bite, a nutrition tracking mobile app, completing full project delivery 1 week ahead of deadline.
    \item Collected, cleaned, and labeled a dataset of 6,500+ food images for multi-class classification model training.
    \item Trained a Computer Vision classification model and deployed it using Flask on Google Cloud Platform (GCP) to serve predictions for the mobile app.
    \item Awarded Best Presenting Group out of 30+ capstone teams for technical delivery and clear project demonstration.
\end{itemize}

\newpage

\cvsection{TECHNICAL PROJECTS}

\textbf{Hybrid Graph-RAG Academic Knowledge Base (Thesis)} \hfill Aug 2025 -- Jul 2026
\begin{itemize}
    \item Built a Hybrid Graph-RAG system combining Neo4j and Milvus to search and connect 10,000+ academic papers and author profiles using both semantic similarity and graph links.
    \item Wrote Apache Airflow DAGs to extract, clean, and load publication data from online academic sources into Neo4j and PostgreSQL.
    \item Co-authored a research methodology paper currently under peer review in a Scopus Q1 journal, documenting graph retrieval benchmarks and entity resolution methods.
    \item Structured data storage using PostgreSQL for metadata, Neo4j for graphs, Milvus for embeddings, and Redis for caching, keeping API response times under 1 second.
\end{itemize}

\vspace{4pt}
\textbf{End-to-End Sentiment Analysis Pipeline for KIP-K Scholarship Tweets} \hfill Jun 2024
\begin{itemize}
    \item Built a Python script and text preprocessing pipeline to collect and clean 4,000+ raw tweets from X.
    \item Trained and evaluated a BiLSTM model using GloVe word embeddings to classify public sentiment into positive, neutral, and negative categories.
    \item Visualized sentiment patterns across 50+ universities in a Looker Studio dashboard to help analyze public feedback on scholarship policies.
\end{itemize}

\cvsection{ORGANIZATIONAL EXPERIENCE}

\textbf{Himpunan Mahasiswa Sains Data (UNESA)} \hfill Surabaya, East Java \\
Staff of Advocacy and Student Welfare Division \hfill Jan 2024 -- Jan 2025
\begin{itemize}
    \item Planned and executed 10+ student events, including data science webinars, workshops, and competitions over a one-year tenure.
    \item Initiated the InSCAW career workshop series, coordinating with industry speakers and engaging 50+ student participants.
    \item Served as a liaison between students and faculty department heads to address academic concerns and welfare requests.
\end{itemize}

\cvsection{SKILLS}

\textbf{Technical Skills:}
\begin{itemize}
    \item \textbf{Programming:} Python, SQL, Bash
    \item \textbf{Data Engineering:} Apache Airflow, PySpark, ETL Pipeline Design, Data Modeling, Kafka
    \item \textbf{Databases \& Storage:} PostgreSQL, BigQuery, Neo4j, Milvus, Redis, MinIO
    \item \textbf{AI \& Machine Learning:} LLMs, RAG / GraphRAG, Computer Vision, NLP, PyTorch, TensorFlow, MLflow
    \item \textbf{Cloud \& DevOps:} AWS, Google Cloud Platform (GCP), Docker, Linux, Git, CI/CD (GitLab CI, GitHub Actions)
    \item \textbf{Soft Skills:} Problem Solving, Cross-Functional Collaboration, Technical Communication
\end{itemize}

\vspace{3pt}
\textbf{Certifications \& Training:}
\begin{itemize}
    \item DeepLearning.AI TensorFlow Developer \& Generative AI Specializations
    \item Google Cloud Skills Boost: Machine Learning Specialization (JuaraGCP)
    \item AWS Certified CloudOps / Cloud Computing Training
    \item Dicoding: Data Analysis with Python
    \item EnglishScore / IELTS Equivalent (Band 7.0 -- C1 Equivalent)
\end{itemize}

\vspace{3pt}
\textbf{Languages:}
\begin{itemize}
    \item Indonesian (Native), English (Fluent, IELTS: 7.0), Javanese (Conversational)
\end{itemize}

\end{document}
